%========================= APPENDIX B: REPOSITORY =============================
\chapter{Repository Layout and Verification Artifacts}
\label{app:repo}

\section{Repository Layout}

\begin{lstlisting}[style=code]
tutor/
|-- README.md
|-- docs/                  # design set (UML): vision, architecture,
|                          # domain model, agent design, sequences,
|                          # API contract, Android architecture, SDLC
|-- backend/               # FastAPI + LangGraph service
|   |-- app/
|   |   |-- api/           # routers, dependencies
|   |   |-- core/          # config (typed Settings), logging, errors
|   |   |-- domain/        # entities, value objects, schemas
|   |   |-- agents/        # LangGraph graph, nodes, prompts
|   |   |-- providers/     # LLM + TTS adapters (mock + Gemini)
|   |   |-- services/      # generation facade, assembler
|   |   |-- repositories/  # lesson cache + asset store
|   |   `-- tools/         # list_models helper
|   |-- tests/             # 48 hermetic, key-free tests
|   |-- Dockerfile
|   `-- docker-compose.yml
|-- android/               # Kotlin + Compose client
|   `-- app/src/main/...   # ui/ (topic, player, history),
|                          # domain/, data/ (remote, local)
`-- .github/workflows/     # backend-ci, backend-deploy,
                           # android-ci, android-release
\end{lstlisting}

\section{Backend Test Coverage Summary}

The 48-test hermetic suite covers, by area:

\begin{itemize}
  \item \textbf{Validator} --- accepts valid SVG/segments; rejects malformed
        XML, missing \texttt{viewBox}, unknown identifiers, empty and
        non-contiguous segments.
  \item \textbf{Generation graph} --- happy path from mock providers;
        repair-loop recovery after an injected broken reference; clean
        failure when the retry budget is exhausted.
  \item \textbf{API contract} --- request/response schemas for all
        endpoints; 200-versus-202 cache behaviour; the error envelope for
        each error code; job lifecycle.
  \item \textbf{Repositories} --- memory and filesystem lesson/asset stores,
        including round-tripping manifests and assets.
  \item \textbf{Gemini adapters} --- structured-output handling, PCM-to-WAV
        wrapping with byte-derived durations, and retry/backoff, exercised
        against a faked SDK client (no network, no keys).
  \item \textbf{Configuration} --- provider/engine/storage selection from
        settings; the session fixture that strips environment variables and
        disables \texttt{.env} loading.
\end{itemize}

\section{Continuous Integration Workflows}

\begin{table}[htbp]
\caption{GitHub Actions workflows}
\label{tab:ci}
\centering
\small
\renewcommand{\arraystretch}{1.3}
\begin{tabular}{@{}lp{0.28\textwidth}p{0.44\textwidth}@{}}
\toprule
\textbf{Workflow} & \textbf{Trigger} & \textbf{Action}\\
\midrule
\texttt{backend-ci} & Push / pull request & Run the hermetic pytest suite;
build the Docker image\\
\texttt{backend-deploy} & Manual / tag & Push image to the container
registry; deploy over SSH to the cloud host (auto-HTTPS proxy)\\
\texttt{android-ci} & Push / pull request & Assemble the debug APK and
upload it as an artifact\\
\texttt{android-release} & Tag & Build a signed release APK and attach it to
a GitHub Release (debug-key fallback)\\
\bottomrule
\end{tabular}
\end{table}
